\part Introduction {The shortest distance between any two points is
always blocked. {\bf Warren Monks, 1988}}

\chap Introduction

The aim of this project is to implement a \nn\ of an unspecified type.
I have chosen to implement the back-propagation model, due to Rumelhart
et al.
The work is intended to make qualitative statements on
several facets of the model.
These are: the model's ability to learn and its dependence upon
parameters, the model's accuracy of operation and its dependence upon
parameters, and finally the model's behaviour upon introduction of
casualties of varying severity into the network.

In this chapter I intend to answer some fundamental questions about
\nns, namely: what their biological relevance is, why and how
they solve problems, different models that have been suggested by other
authors, and how these networks can be simulated in hardware and
software.

In the next chapter I will discuss the back-propagation algorithm in
detail, and the chapter following will contain a detailed implementation
independent software specification for a software model.

\sec Definition of Terms

Before I begin this discussion of \nns, definition of a few
terms will aid clarity~[1].

A {\sl connection\/} is a signal pathway between processing elements (or
\neus), that correspond to the axons and synapses of biological
\neus.
Many connections form a \nn.

A {\sl processing element\/} is a simple artificially simulated \neu\ %
that has a graded (analog) response.
(From this point forward, I will draw a biological analogy and refer to
processing elements as \neus.)
It consists of a memory for its present state, and a transfer function
that maps some function of the inputs to the output.
The output is connected to many other processing elements' (or \neus')
inputs.

A {\sl weight\/} is a dynamic value (at least during learning) that
determines the intensity of connections between \neus.
Weights can be positive or negative (or zero,) indicating an excitatory or
inhibitory (or no) connection.
A weight is associated with a single input connection to a \neu.

\sec So, What Is a \NN?

A \nn\ is a collection of graded response \neus\ (processing elements),
representing an approximation to biological \neus, in some kind of
feedforward/feedback network.

There are many different kinds of \nns\ (Hecht-Nielsen~[1] suggests 
that there are 13
main types), with the most common (or important) being the Hopfield,
back-propagation (Rumelhart et al.), ``improved'' neocognitron (Fukushima), and
adaptive resonance theorem or ART (Carpenter and Grossberg).
Each is particularly good at a number of tasks, and some have weak
points.
They vary considerably in complexity and power.

\Nn s are often described as a network of ``collective decision''~[2]
circuits.
They gain their power the same way the human brain does: many (many
many) simple analog devices computing at the same time, while connected
in parallel and continuously communicating with each other.
The human brain has, by most estimates, $10^{11}$ \neus\ %
($100$~billion!) with up to~$10\,000$ interconnections per \neu.
Imitating the performance of a human ``biological \nn'' with an
artificial one is not and may never be feasible.
However, some of the characteristics of the brain, such as the ability
to learn or be taught (there is a difference here, as some networks
teach themselves and some require a teacher) can be achieved.
The brain often returns several possible solutions to a problem, with
varying degrees of certainty.
Artificial \nns\ can do this.
And both man and \nns\ can infer the original from an image that is
noisy, or showing only parts of the whole.
\goodbreak\midinsert\vskip 6cm
Figure 1. Classic neural circuit.\hfill Figure 2. I/O response of simulated
\neus.
\endinsert
Figure~1 shows the classic feedforward and feedback neural circuit, and
figure~2 demonstrates the sigmoid output response of a simulated \neu.
These diagrams are referred to in a later section, ``Why and how do
\nns\ solve problems?''

Simulations of \nns\ involve several important simplifications to the
behaviour of the individual \neus.
It has been shown that analog \neus\ perform better than earlier 
two-state (binary \neu) attempts~[2,17], but many features evident in 
biological \neus\ %
(much chemistry is discussed in~[3])
are simply ignored in the most commonly used models of \neus, and analogies
are made to voltages and currents.
The model I will use for \neus\ is simple also, with the transfer
function shown in figure~2 being the output from the inputs times the
weights and summed.
The relationship of the simple model to biology is shown in~[4].

\sec Neurological Relevance

Neurons themselves, in most models, are extremely simplified versions of
the real biological article.
So, it is the overall structure of a \nn\ that makes it relevant to
neurology, not its individual components.
This has been proven by many people; examples will be described
below.
Techniques for using \nns\ in neurology take two approaches: the structure of the
brain is discovered by experiment, and a network is built to verify this
(there is suprisingly good correlation between the results of
experiments such as this and the network), or a network of the
appropriate model is built and taught and then the parts of the brain
corresponding to the model are easily identified.
Both methods, by dealing with small parts of the brain at a time, have
yielded very good results.

Parts of the brain, specifically area~7a of the posterior parietal
cortex in monkeys, have been successfully modelled by Zipser and
Andersen~[5] using
back-propagation [20] learning in a feed-forward \nn, similar in form to the
one intended for this project.
The neurons respond to the location of the stimulus with respect to the
eye and the position of the eyes to calculate the location of external
objects.
Even though the authors admit that there is no way
that the back-propagation method is the only method used in the brain,
it is very likely that a combination of Hebb-like~[6,18] learning and
the feeding back of errors would generate similar results. (Hebb was the
first to suggest a biologically plausible method of learning. 
It is described in a later section.)
They note also that all cortical connections in the brain have
reciprocal feedback pathways, presumably for the feeding back of error
signals.

Bear, Cooper and Ebner~[3] have used \nns\ in the reverse direction to
study the primary visual cortex area 17 in adult cats.
They used a \nn, determined theoretically how the \neus\ could behave,
and compared this with the actual data.
The theoretical model, then, enabled the researchers to sort out which
of the possible hypotheses of brain function was most correct by experiment.

Frohn, Geiger and Singer~[7] have used a model based not on rigorous
maths, but, rather, neurological data, to account for the features in a
mammal's visual system.
They arrived at a~5 layer model, trained by a Hebb-like rule~[6,18] that
teaches itself.
This allowed the authors to conclude that ``the internal representation
of a stimulus is the spatial activity pattern of a reciprocally coupled
population of \neus,'' i.e.\ the knowledge of patterns in the visual system
is distributed.
Sun, Chen and Lee [19] also give a mathematical description of Hebbian
learning of stereopsis in a \nn.

The importance of motion to the perception of structure in mammals has
been simulated~[8] with a three-layer back-propagation \nn, and the
performance of the model was almost identical to that of man and monkey.
The first and third layers of the model are designed to simulate 
\neus\ in a particular area of the brain, while it has not been accurately
determined where the middle layer could lie. 

Some models of \nns, such as ART (due largely to Grossberg), are strongly 
based on structures found in certain parts of the brain. 
In fact, experiments with ART networks have led to a number of important
predictions~[9] of human brain structures that have been proven correct.
Needless to say ART is an extremely complex form of network.

Other authors such as Lyon and Mead [21] have taken an entirely different
approach.
They argue that to build machines that operate like humans then the
machines will need to ``perceive'' like humans.
Hence they have built an electronic cochlea based around analog
neuron-like components which operate like the human ear.
This approach appears very fruitful, and I suspect we will see more
machines designed about human form (cybernetics) in the future.

\sec Why and How do \NN s Solve Problems?

As mentioned earlier, \nns\ gain their power from the connection of
simple analog devices in parallel.
The human brain has roughly~$10^{11}$ \neus, and each \neu\ %
can have up to~$10\,000$ interconnections.
Each connection can be an excitatory or inhibitory input or feedback, or
an output connection to other \neus.

At present, it seems the largest neurocomputer\footnote*{Dedicated
machines built to run models of neural networks appear to have acquired
the jargonistic term ``neurocomputers''.} built to date has~$10^6$
\neus\ and interconnections totalling~$1.5\times10^6$ in total~[1].
This is, of course, far from the complexity of the human brain, but as
we shall see it still offers a great deal of power in solving problems.

There are lots of common problems which cannot, as yet, be solved  by
standard serial algorithmic means (some suggest they can never be).
For instance, recognizing the difference between friendly and enemy
aircraft, a spoken language translator, or a system recognizing speech
no matter who is speaking.

Human beings perform these tasks without any apparent difficulty, once
they have learned to do so.
Often, for many of these actions, there is no need to be explicitly taught.
For all of these problems, we as man or woman can generate many examples 
of what the
output should be for a given input, even though we cannot write down a
fixed algorithm on how to do it.

But \nns, without being programmed with an algorithm but taught instead,
can already perform many of these tasks with~$95\%$ accuracy. For
example, Kohonen~[10] has built a phonetic typewriter, Fukushima~[11]
has used his ``improved''\footnote*{In the latest writings I possess~[11],
Fukushima refers to the model he is working on as an improved version of
his slightly older neocognitron model (circa 1984). Hence I have dubbed it an
``improved'' neocognitron, as no other official term was used in
[11].} 
neocognitron model to recognise Chinese characters, and many
authors have shown the ability of \nns\ to remove noise and recognize
patterns~[7,11,12,23,24,25].
Many other implementations of this new non-algorithmic \nn\ programming are
being carried out.
A very good example is using \nns\ as a signal analyser to detect
subtle patterns in neuroelectric signals from brian data with ``highly
dimensioned noise'' [26].

To understand how \nns\ solve problems, we'll first try to get an
intuitive feel for how they work, compared to a digital or Von~Neuman
machine that we are all, no doubt, familiar with already.

A digital machine takes an input, such as programs and data, and goes
through a rigidly defined set of steps (instructions), branching if
necessary, and producing an output of some kind, based upon program and
data.
On the other hand, \nns\ or collective decision circuits, take 
inputs and move to minimize
some function of the ``computational energy'' of the inputs, based upon
the dynamic weights of the connections between the \neus. 
Some authors, notably Hopfield and Tank~[2]\footnote{**}{and recognized 
in~[13] as relating to other physical systems as well, notably spin glasses},
see \nns\ as an energy
reducing machine: imagine an $n$-dimensional surface of hills and
valleys, with the computation beginning
at some high point on the surface.
As the computation procedes it will move to a valley, downhill (i.e.\ %
minimizing ``computational energy''), until it is stable.
The final answer then appears at the output \neus.

Consider figure~1 again.
It contains several types of \neus, namely what we can call
``principal'' and ``inhibitory'' \neus\ (labelled by `P' and `IN'
respectively).
Looking at figure~2, imagine the response of the `P' \neus\ follows the
solid line, and the response of the `IN' \neus\ follows the dotted line.
Then, the response of the circuit to an input will depend upon the
weights on the connections between \neus.
As you can see, the deviation from zero of an input to any 
\neu\ (and the input is
the sum of the outputs of the previous \neus\ times their connection
weights) will tend to drive it strongly towards a positive or negative
output.
All the \neus\ ``calculate'' at the same time, so the behaviour of a
neural system (especially a very large one) is not trivial, and can't be
written as a serial algorithm.

\sec Types of \NN s

There are many different types of \nns.
Some are particularly suited to some types of problem, some are suitable
for many types of problem.
Some can ``learn'' without supervision, others require a teacher.
Most of these neural models are grounded upon Donald
Hebb's~[6] theory of synapse modification.
He indicated that a plausible explanation for learning would be for
connections between \neus\ to strengthen if the activity of both ends of
the connection increased (ie., reinforce the connection) and weaken the
connection if the ends weaken in activity (ie., inhibit the connection).
To various extents, all these models aim to implement ``Hebbian''
learning by one method or another.
(Hebb never mathematically quantified his thinking.
A good approach to doing just this is given by Linsker [18].)
Examples of several types of \nns\ appear below.

\ssec The Hopfield Network

The Hopfield network shown in figure~3 is the simplest type of network,
yet it is suprisingly powerful.
\figure {3} {5} {Hopfield network}
It has a feedback structure, and it cannot learn (although a technique
for teaching feedback networks has been suggested by Atiya~[14]).
The selection of weights depends upon the problem at hand, for example
classic computer science problems such as the Travelling Salesman
Problem (TSP), the book stacking problem, and a symmetry detector.
All these problems can be computed after the appropriate weights are set
up in the feedback structure.

Consider the Travelling Salesman Problem, or TSP.
A salesman is required to visit a number of cities, by the shortest
possible route.
Imagine that there are 30 cities on the tour.
Comparing the \nn\ solution to the serial solution shows some staggering
results!

Let's see the difference between a brute force approach, a good serial
algorithm, and a Hopfield \nn\ solution.
Based upon estimates from~[15],
probably~$10^{30}$ comparisons would be needed for a brute force
approach.
This is obviously not computable.
An excellent serial algorithmic estimation approach would 
require roughly~$1\,500$
comparisons in the best case (but probably many more).
But a Hopfield \nn, it has been found, will select one of the
best~$10^7$ solutions (and this is as good as or better than the algorithmic
approach described above) in a single convergence of the network, or
several time-constants of the ciruit of figure 3.
Quite a difference.

\ssec The Back-Propagation Network

This is another ``classic'' \nn, along with the one described above
(they appear to be the most commonly implemented in neurological
experiments).
The inventors and main developers~[1] are considered to be Werbos,
Parker and Rumelhart.
It has a feedforward structure, as shown in figure 4.
\figure {4} {5} {A typical back-propagtion network}
This network can be taught [20], ie strengths of connections may be varied
within a rigidly defined set of rules called a ``learning algorithm'',
rather than being set by hand.
The next chapter explains the back-propagation algorithm in detail.

The network learns by comparing the network output with the desired
output, for a fixed input.
The error is used to change the weights, then another input/output pair
is presented.
The process continues through the samples as many times as necessary
until the error falls below a certain threshold.
How many times each input/output pair needs to be presented is a matter
of conjecture, and one of the things I'm going to investigate later in
the project.

By itself, back-propagation networks are not neurologically plausible.
However, the propagation of errors back through the network by feedback
is known to occur in the cerebal cortex.

Like all other \nns, this one is redundant (ie knowledge is distributed)
and so behaves sanely when
casualties in processing elements occur.
Once again, exactly how robust a network is is for me to investigate later
in the project.

Back-propagation networks are good for solving problems such as speech
synthesis and adaptive control.

\ssec The ``Improved'' Neocognitron

The ``improved'' neocognitron, along with the next model discussed, falls 
into that category of \nns\ that are more neurologically sound in their
design.
They both contain more complex \neus, with complex feedforward and
feedback connections.

Fukushima~[11] has used his model to recognise Chinese characters, a
task that is readily performed by over~1~billion Chinese every day but
was yet to be done reliably by machine.
It is largely immune to those bugbears of artificial \nns, namely
rotation, translation, and changes in scale.
One of the most astounding properties of this model is its ability to
shift ``attention'', just like a human, from one pattern to another for
an input that contains several recognisable patterns at the same time.
(See~[11], which has examples.)
\figure {5} {6} {An example of an ``improved'' neocognitron}

Figure 5 shows us the nature of the network: feedback and feedforward
paths, with dynamic gain and threshold controls forming part.
There are fixed and dynamic connections between \neus.

The model functions as follows: As patterns are recognized by the
network (by itself ie., without a teacher), the top or ``recognition''
layer of \neus\ becomes active.
The output of this layer is fed back down to the lower stages, which, by
varying the gain, enforce the forward flow of the recognized pattern.
After the network has learnt to recognize the presented patterns, by
simply breaking the feedback path for an instant the ``attention'' is
switched to another recognizable pattern, and a new pattern will be
recognized at the output.
Needless to say, this network is too complex to implement in the short
time available, so this is as thoroughly as it will be discussed.

\ssec The Adaptive Resonance Theory Model

The Adaptive Resonance Theory (or ART) model is the most complex one
discussed here~[9].

ART systems belong to a group of \nns\ which fall into a category of
``competitive learning'' models.
It has been used for visual pattern recognition, speech perception, and
radar classification.

Even though it is the most complex model examined here, it can be
simply described as a two layer network with gain control.
There is a distinction between ``short'' and ``long'' term memory, with
memory decaying slowly over time (like a human brain, if some things are not
used all the time they are gradually forgotten\dots)
There is extra circuitry as well for resetting the second layer (ie
disabling its reinforcing action) when there are sufficiently large
mismatches between the two levels.

\sec Simulating \NN s in Hardware

As explained earlier, \neus\ have a graded response.
Hence it is possible to simulate a \neu\ with a transfer function like
that of figure~2 as an amplifier with a similar response.
It would be possible, if one stretched one's imagination a little, to
build a huge multi-million \neu\ ``brain'' from op-amps and resistors,
similar in structure to the network shown in figure 1.
However, this is not a particularly sensible aim, for instance how do
the weights change automatically? (Who twiddles the knobs on the
variable resistor as the network learns?)
For {\sl small\/} problems, this technique could be used.
\figure {6} {5} {A special machine for the TSP}
For example Hopfield and Tank~[2] devised such a contraption shown in
figure~6 using a Hopfield network that could solve the TSP for however
many cities it was wired for (I don't think they actually built it).
The output is indicated by a globe lit for the appropriate city (column)
in the appropriate order (row) of the salesman's trip.
The device's weights are set so that only one globe can light fully in
any row or column.
Needless to say after that example the discrete approach to \nns\ is not
particularly popular.

There have been some attempts at analog ICs with variable internal
weights (one such is described in~[2]).
The ones built have functioned successfully as associative memorys.
However the density and accuracy of VLSI analog ICs is low, so
this approach has not been terribly fruitful.

A novel approach to hardware simulation has been described by Murray and
Smith [16], whereby they use digital ICs which use pulse stream
arithmetic.
Streams of pulses of varying frequency are gated within the connections.
This gives a situation near to biology, where a \neu\ that is {\sl ``on''\/}
produces a regular train of digital pulses (the frequency dependent upon
its {\sl ``on''}ness), while a \neu\ that is {\sl ``off''\/} produces nothing.
This technique is still in the experimental stage at the date of
publication, but it could prove possible to efficiently perform
significant calculations in this way in hardware.

Another novel approach is described by Vidal~[22].
He proposes a \nn\ implemented in programmable logic, using purely
digital techniques, and argues that possibly digital (Boolean)
techniques offer a good solution to network problems.
This is an example of how varied the models for both hardware and
software are, and how much there is that is still undiscovered and
undecided in this field.

The perfect hardware for the parallel operation of a \nn\ is not planar
silicon.
Most authors suggest that three dimensional biological materials are
infinitely better suited to parallel processing than a two dimensional
substance.
However optical computers, which are parallel by their very nature, could
offer processing solutions in the long term.

It is now easy to understand why most \nns\ are simulated in software.
Most of the references given refer to software models as they are more
flexible, though slower.
The next section deals briefly with this.

\sec Simulating \NN s in Software

As I stated at the start of this chapter, this project is intended to
produce a working neural network of the back-propagation model~[20].
Simulation of the parallel activity of the \neus\ is possible by 
imagining the operations
occuring in discrete time steps, and calculating one row of \neus\ at a
time.
Of course, this technique works best when signals flow in only one direction,
such as in back-propagation models, because then one can calculate the
first layer, then the values from the first give the second, etc.

The simulation of other networks such as the ``improved'' neocognitron
will not be considered here, but I suspect it would not be much more
difficult: possibly requiring calculations of layers with smaller time
steps due to the more complex interactions between layers.

The next chapter describes the mathematics behind the back-propagation
algorithm to be implemented, based largely on~[20] by Rumelhart, Hinton
and Williams.
